% !TeX root = ../教材题库_全部条目_转写工作稿.tex
% Source: XJ-B1-C04-EX-P130-Y1358
% 第4章 / 4.4 函数与方程 / 教材 p.130 / PDF 第137页
% OCR未校对：公式与文字均可能有误；未完成数学TeX转写；禁止直接发给学生
% 原题及相邻原解截图（核对用）：../题目原图/XJ-B1-C04-EX-P130-Y1358-01.webp
% 本文件只在转写完成后作为题干使用；OCR 条目尚不是可靠题干。
\begin{bankitem}{4.4 例题 例题1}{第4章 / 4.4 函数与方程 / 教材 p.130 / PDF 第137页}{XJ-B1-C04-EX-P130-Y1358}{OCR未校对：公式与文字均可能有误；未完成数学TeX转写；禁止直接发给学生}
\begin{ocrdraft}
[教材p.130 / PDF第137页]
\par
0〔〔]讨论函数丆）一2一5在区间（1，2）内零点的个数．
\par
解由于不1)一一3<0，不2）一11>0，又丿、0）一2於一5是单调递增函数，因
\par
此，函数不x）在区间（1，2）内只有一个零点．
\par
运用函数的思想方法来求方程的解可以给我们带来很大的便利．例如，方程
\par
丆0）一g(x）的解既可以看作两个函数3-不和3，一g(x）的图象的公共点的横坐标，
\par
也可看作函数3，一丆0）一g(x）的零点．
\end{ocrdraft}
\end{bankitem}
